\subsection{Integrable sets}

DEFINITION 2. --- A subset A of a locally compact space X is said to be integrable with respect to a measure \(\mu\) on X (or, that it is \(\mu\) -integrable) if the characteristic function \(\varphi_{A}\) of A is integrable. The finite number \(\mu(A) = \int \varphi_{A} d\mu\) is called the measure of A.

For every integrable set A, \(|\mu|(A) = |\mu|^{*}(A)\) (No.~2, Prop. 1); for a set to be negligible, it is necessary and sufficient that it be of measure zero with respect to \(|\mu|\) .

PROPOSITION 6. --- The union of a finite family \((\mathrm{A}_{i})_{1\leqslant i\leqslant n}\) of integrable sets is integrable, and

\[
| \mu | \left(\bigcup_ {i = 1} ^ {n} \mathrm{A} _ {i}\right) \leqslant \sum_ {i = 1} ^ {n} | \mu | (\mathrm{A} _ {i}).\tag{13}
\]

If, moreover, the \(\mathrm{A}_i\) are pairwise disjoint, then

\[
\mu \left(\bigcup_ {i = 1} ^ {n} \mathrm{A} _ {i}\right) = \sum_ {i = 1} ^ {n} \mu (\mathrm{A} _ {i}).\tag{14}
\]

For, if \(\mathbf{A} = \bigcup_{i=1}^{n} \mathbf{A}_i\) then \(\varphi_{\mathbf{A}} = \sup \varphi_{\mathbf{A}_i}\) , therefore (§3, No.~5, Cor. of Prop. 12) if the \(\mathbf{A}_i\) are integrable then so is \(\mathbf{A}\) ; the relation (13) is a special case of the analogous relation for outer measures (§1, No.~4, Prop. 18), on taking account of the relation \(|\mu|(\mathrm{A}) = |\mu|^*(\mathrm{A})\) ; finally, if the \(\mathbf{A}_i\) are pairwise disjoint, then \(\varphi_{\mathrm{A}} = \sum_{i=1}^{n} \varphi_{\mathrm{A}_i}\) , whence (14).

PROPOSITION 7. --- \(1^{\circ}\) If A and B are two integrable sets such that \(\mathbf{B} \subset \mathbf{A}\) , then the set \(\mathbf{C} = \mathbf{A} - \mathbf{B}\) is integrable and

\[
\mu (\mathrm{C}) = \mu (\mathrm{A}) - \mu (\mathrm{B}).\tag{15}
\]

\(2^{\circ}\) The intersection of a countable family of integrable sets is integrable. The first part follows from the fact that \(\varphi_{\mathrm{C}} = \varphi_{\mathrm{A}} - \varphi_{\mathrm{B}}\) . On the other hand, if \((\mathbf{A}_n)\) is a sequence of integrable sets and A is its intersection, then \(\varphi_{\mathrm{A}} = \inf_{n} \varphi_{\mathrm{A}_n}\) , therefore A is integrable (No.~3, Prop. 4).

COROLLARY. --- If \((\mathrm{A}_{n})\) is a decreasing sequence of integrable sets, then \(\mu(\bigcap_{n}\mathrm{A}_{n})=\lim_{n\to\infty}\mu(\mathrm{A}_{n})\) .

For, if \(\mathrm{A} = \bigcap_{n}\mathrm{A}_{n}\) then \(\varphi_{\mathrm{A}}\) is the lower envelope of the decreasing sequence \((\varphi_{\mathrm{A}_n})\) (No.~3, Prop. 4).

PROPOSITION 8. --- Let \((\mathbf{A}_{n})\) be an increasing sequence of integrable sets; for the union \(A = \bigcup_{n} A_{n}\) to be integrable, it is necessary and sufficient that \(\sup_{n} |\mu|(\mathbf{A}_{n}) < +\infty\) , in which case,

\[
\mu (\mathrm{A}) = \lim _ {n \rightarrow \infty} \mu (\mathrm{A} _ {n}).\tag{16}
\]

For, the \(\varphi_{\mathrm{A}_n}\) form an increasing sequence of integrable functions, and \(\varphi_{\mathrm{A}} = \sup_n \varphi_{\mathrm{A}_n}\) ; thus, the proposition follows from Prop. 4 of No.~3.

COROLLARY. --- Let \((\mathrm{A}_{n})\) be a sequence of integrable sets such that \(\sum_{n=1}^{\infty} |\mu|(\mathrm{A}_{n}) < +\infty\) ; the union \(A = \bigcup_{n} A_{n}\) is integrable, and

\[
| \mu | \left(\bigcup_ {n} \mathrm{A} _ {n}\right) \leqslant \sum_ {n = 1} ^ {\infty} | \mu | (\mathrm{A} _ {n}).\tag{17}
\]

For, \(\varphi_{\mathrm{A}} = \sup_n\varphi_{\mathrm{A}_n}\) and

\[
| \mu | ^ {*} (\mathrm{A}) \leqslant \sum_ {n = 1} ^ {\infty} | \mu | ^ {*} (\mathrm{A} _ {n}) = \sum_ {n = 1} ^ {\infty} | \mu | (\mathrm{A} _ {n}) <   + \infty
\]

(§1, No.~4, Prop. 18); therefore A is integrable (§3, No.~6, Cor. 2 of Th. 5) and, since \(|\mu|(\mathrm{A}) = |\mu|^*(\mathrm{A})\) , we indeed have (17).

PROPOSITION 9. --- Let \((\mathbf{A}_{n})\) be a sequence of pairwise disjoint integrable sets such that \(\sum_{n=1}^{\infty}|\mu|(\mathbf{A}_{n})<+\infty\) ; then

\[
\mu \left(\bigcup_ {n} \mathrm{A} _ {n}\right) = \sum_ {n = 1} ^ {\infty} \mu (\mathrm{A} _ {n}).\tag{18}
\]

For, if \(\mathbf{A} = \bigcup_{n}\mathbf{A}_{n}\) then \(\varphi_{\mathrm{A}} = \sum_{n = 1}^{\infty}\varphi_{\mathrm{A}_n}\) , and the proposition follows from (17) and Cor. 2 of Th. 2 of No.~3.

The relation (18) is also expressed by saying that the measure \(\mu\) is completely additive in the set of integrable subsets of X.

